% Generated draft: requires mathematical and source review.
\documentclass[11pt]{article}
\usepackage[margin=25mm]{geometry}
\usepackage{fontspec}
\setmainfont{Latin Modern Roman}
\usepackage{amsmath,amssymb,mathtools,unicode-math}
\setmathfont{Latin Modern Math}
\usepackage{enumitem,needspace,xurl,hyperref}
\setlength{\parindent}{0pt}
\setlength{\parskip}{5pt}
\setlength{\emergencystretch}{3em}
\newcommand{\N}{\mathbb{N}}
\newcommand{\Z}{\mathbb{Z}}
\newcommand{\Q}{\mathbb{Q}}
\newcommand{\R}{\mathbb{R}}
\newcommand{\C}{\mathbb{C}}
\newcommand{\F}{\mathbb{F}}
\newcommand{\A}{\mathbb{A}}
\newcommand{\PP}{\mathbb{P}}
\newcommand{\E}{\mathbb{E}}
\newcommand{\eps}{\varepsilon}
\newcommand{\dd}{\,\mathrm{d}}
\begin{document}
\section*{UnsolvedMath ID: 700854}
\textbf{Source problem code:} catalogue entry
\subsection{Definitions and mathematical statement}
A formation of finite groups is a nonempty class \(\mathfrak F\) of finite groups, closed under isomorphisms, homomorphic images, and finite subdirect products. Here a subdirect product is a subgroup of a direct product whose projection onto each factor is surjective. For a finite group \(X\), define its \(\mathfrak F\)-residual by
\[
X^{\mathfrak F}=\bigcap\{N\trianglelefteq X:X/N\in\mathfrak F\}.
\]
A subgroup \(H\leq G\) is subnormal in \(G\) if there is a finite chain
\[
H=H_0\trianglelefteq H_1\trianglelefteq\cdots\trianglelefteq H_r=G,
\qquad r\geq 0.
\]
For subgroups \(H,K\leq G\), write \(\langle H,K\rangle\) for the subgroup of \(G\) generated by \(H\cup K\).

\textbf{Source qualification.} The source requires \(\mathfrak F\) to be a \emph{radical composition formation} of finite groups. The supplied record does not define the additional technical qualifications ``radical'' and ``composition''. They are retained as hypotheses in the source's intended sense; their precise definitions require confirmation before this becomes a fully self-contained statement.

\textbf{Requested assertion (S. F. Kamornikov, Kourovka Notebook Problem 12.35).} For every radical composition formation \(\mathfrak F\) of finite groups, every finite group \(G\), and every pair of subnormal subgroups \(H,K\) of \(G\), prove that
\[
\bigl\langle H,K\bigr\rangle^{\mathfrak F}
=
\bigl\langle H^{\mathfrak F},K^{\mathfrak F}\bigr\rangle.
\]
Each residual is computed in the group appearing as its base: respectively \(\langle H,K\rangle\), \(H\), and \(K\). The equality is an equality of subgroups of \(G\).
\subsection{Short English statement}
For any radical composition formation \(\mathfrak F\), the \(\mathfrak F\)-residual of the subgroup generated by two subnormal subgroups of a finite group should equal the subgroup generated by their individual \(\mathfrak F\)-residuals. The source does not define the technical qualifications on the formation, so their intended meaning requires review.
\subsection{Sources}
\begin{itemize}
\item Supplied local UnsolvedMath registry record: permanent ID \(700854\), exact source code \texttt{catalogue entry}, title \emph{Kourovka Notebook Problem 12.35}; upstream revision \texttt{d371ee917ac9647fbb06cd6fc529deec6f3dda82}.
\item \emph{The Kourovka Notebook (Unsolved Problems in Group Theory)}, original problem 12.35, 12th issue (1992); proposed by S. F. Kamornikov. The supplied record also identifies the 21st issue edition (2026), arXiv identifier \texttt{1401.0300v48}. Its collection description names E. I. Khukhro and V. D. Mazurov as editors.
\item The local record labels the problem \emph{open} and records an update dated 2026-10-16. These are supplied registry metadata, not an independently verified mathematical or literature assessment. No separate literature-verification date or reference defining ``radical composition formation'' is supplied.
\end{itemize}
\end{document}
